% GENERATED FILE. Edit canonical lessons, then run npm run generate:books.
% canonical-source-hash: fnv1a64:9fb18eda4fffcd73
% canonical-lessons: ZH-C179-xiyiji, ZH-C179-xiyifen, ZH-C179-liang, ZH-C179-yun, ZH-C179-yigui

\chapter{Laundry Day}
\label{ch:zh-laundry-day}
\hlchaptermodality{179}

\begin{chapteropening}
\textbf{By the end of this chapter:} \emph{I can talk about the washing machine and powder, hanging things out, ironing and the wardrobe.}

Complete the last lesson of chapter 179: I can talk about the washing machine and powder, hanging things out, ironing and the wardrobe.
\end{chapteropening}

\section[\texorpdfstring{\zh{洗衣机} (xǐyījī)}{洗衣机 (xǐyījī)}]{\zh{洗衣机} (xǐyījī) — washing machine}

\label{lesson:ZH-C179-xiyiji}



\begin{quote}
Before the new one: say the Mandarin for lately, then the Mandarin for in the small hours.
\end{quote}

\subsection*{You'll want to know: \zh{洗衣机}}
\textbf{\zh{洗衣机}} — \emph{xǐyījī} — \textquotedblleft{}washing machine\textquotedblright{}.

\begin{quote}
\textbf{\zh{我家有洗衣机}} \emph{wǒ jiā yǒu xǐyījī} — We've got a washing machine at home.
\end{quote}

\begin{cousinweb}
\textbf{\zh{洗衣机}} (\emph{xǐyījī}) is built from \zh{洗} 'to wash', \zh{衣} 'clothes' and \zh{机} 'machine', all three inherited from Middle Chinese, \emph{sejX}, \emph{'j+j} and \emph{kj+j} (Guangyun, in Baxter's transcription); the dictionary glosses the last as the trigger of a crossbow, and 'machine' is a later sense.
\end{cousinweb}

\begin{grammarlens}[title={What you've built}]
One more word for saying what I do.
\end{grammarlens}

\subsection*{Your turn}
\begin{itemize}
\raggedright
  \item \emph{Say it:} \emph{xǐyījī}
  \item \emph{Say it:} \emph{xǐyījī}, once more
  \item \emph{Say it:} \emph{wǒ jiā yǒu xǐyījī} — We've got a washing machine at home.
  \item \emph{Recall:} say the Mandarin for lately, then the Mandarin for in the small hours, then say \emph{xǐyījī} again
\end{itemize}

\begin{tcolorbox}[breakable,colback=teal!4,colframe=teal!35!black,title={Before you move on}]
What does \zh{洗衣机} mean? (\textquotedblleft{}Washing machine\textquotedblright{}.) Say it once more.
\end{tcolorbox}
\section[xǐyīfěn]{xǐyīfěn — washing powder}

\label{lesson:ZH-C179-xiyifen}



\begin{quote}
Before the new one: say the Mandarin for in the small hours, then the Mandarin for a washing machine.
\end{quote}

\subsection*{You'll want to know: xǐyīfěn}
\emph{xǐyīfěn} — \textquotedblleft{}washing powder\textquotedblright{}.

\begin{quote}
\emph{xǐyīfěn méiyǒu le} — We're out of washing powder.
\end{quote}

Its characters come later; for now it is a word for the ear.

\begin{cousinweb}
\emph{xǐyīfěn} is built from \zh{洗} 'to wash', \zh{衣} 'clothes' and \emph{fěn} 'powder', all three inherited from Middle Chinese, \emph{sejX}, \emph{'j+j} and \emph{pjunX} (Guangyun, in Baxter's transcription), whose example for the last is a white powder made from lead.
\end{cousinweb}

\begin{grammarlens}[title={What you've built}]
One more word for saying what I do.
\end{grammarlens}

\subsection*{Your turn}
\begin{itemize}
\raggedright
  \item \emph{Say it:} \emph{xǐyīfěn}
  \item \emph{Say it:} \emph{xǐyīfěn}, once more
  \item \emph{Say it:} \emph{xǐyīfěn méiyǒu le} — We're out of washing powder.
  \item \emph{Recall:} say the Mandarin for in the small hours, then the Mandarin for a washing machine, then say \emph{xǐyīfěn} again
\end{itemize}

\begin{tcolorbox}[breakable,colback=teal!4,colframe=teal!35!black,title={Before you move on}]
What does xǐyīfěn mean? (\textquotedblleft{}Washing powder\textquotedblright{}.) Say it once more.
\end{tcolorbox}
\section[liàng]{liàng — to hang out to dry}

\label{lesson:ZH-C179-liang}



\begin{quote}
Before the new one: say the Mandarin for a washing machine, then the Mandarin for washing powder.
\end{quote}

\subsection*{You'll want to know: liàng}
\emph{liàng} — \textquotedblleft{}to hang out to dry\textquotedblright{}.

\begin{quote}
\emph{jīntiān tiānqì hěn hǎo nǐ qù liàng yīfu} — The weather's nice today: go and hang the washing out.
\end{quote}

Its characters come later; for now it is a word for the ear.

\begin{cousinweb}
\emph{liàng} No dictionary we use records an older form, so its history is not settled.
\end{cousinweb}

\begin{grammarlens}[title={What you've built}]
One more word for saying what I do.
\end{grammarlens}

\subsection*{Your turn}
\begin{itemize}
\raggedright
  \item \emph{Say it:} \emph{liàng}
  \item \emph{Say it:} \emph{liàng}, once more
  \item \emph{Say it:} \emph{jīntiān tiānqì hěn hǎo nǐ qù liàng yīfu} — The weather's nice today: go and hang the washing out.
  \item \emph{Recall:} say the Mandarin for a washing machine, then the Mandarin for washing powder, then say \emph{liàng} again
\end{itemize}

\begin{tcolorbox}[breakable,colback=teal!4,colframe=teal!35!black,title={Before you move on}]
What does liàng mean? (\textquotedblleft{}To hang out to dry\textquotedblright{}.) Say it once more.
\end{tcolorbox}
\section[yùn]{yùn — to iron}

\label{lesson:ZH-C179-yun}



\begin{quote}
Before the new one: say the Mandarin for washing powder, then the Mandarin for to hang out to dry.
\end{quote}

\subsection*{You'll want to know: yùn}
\emph{yùn} — \textquotedblleft{}to iron\textquotedblright{}.

\begin{quote}
\emph{wǒ yào yùn yīfu} — I'm going to do the ironing.
\end{quote}

Its characters come later; for now it is a word for the ear.

\begin{cousinweb}
\emph{yùn} is inherited from Middle Chinese \emph{'jut} (Guangyun, in Baxter's transcription), glossed 'to smooth silk with fire'. That reading would regularly give yù, so how today's yùn came about is not settled.
\end{cousinweb}

\begin{grammarlens}[title={What you've built}]
One more word for saying what I do.
\end{grammarlens}

\subsection*{Your turn}
\begin{itemize}
\raggedright
  \item \emph{Say it:} \emph{yùn}
  \item \emph{Say it:} \emph{yùn}, once more
  \item \emph{Say it:} \emph{wǒ yào yùn yīfu} — I'm going to do the ironing.
  \item \emph{Recall:} say the Mandarin for washing powder, then the Mandarin for to hang out to dry, then say \emph{yùn} again
\end{itemize}

\begin{tcolorbox}[breakable,colback=teal!4,colframe=teal!35!black,title={Before you move on}]
What does yùn mean? (\textquotedblleft{}To iron\textquotedblright{}.) Say it once more.
\end{tcolorbox}
\section[yīguì]{yīguì — wardrobe}

\label{lesson:ZH-C179-yigui}



\begin{quote}
Before the new one: say the Mandarin for to hang out to dry, then the Mandarin for to iron.
\end{quote}

\subsection*{You'll want to know: yīguì}
\emph{yīguì} — \textquotedblleft{}wardrobe\textquotedblright{}.

\begin{quote}
\emph{nǐ de dàyī zài yīguì lǐ} — Your coat is in the wardrobe.
\end{quote}

Its characters come later; for now it is a word for the ear.

\begin{cousinweb}
\emph{yīguì} is built from \zh{衣} 'clothes' and \emph{guì} 'a chest', both inherited from Middle Chinese, \emph{'j+j} and \emph{gwijH} (Guangyun, in Baxter's transcription).
\end{cousinweb}

\begin{grammarlens}[title={What you've built}]
One more word for saying what I do.
\end{grammarlens}

\subsection*{Your turn}
\begin{itemize}
\raggedright
  \item \emph{Say it:} \emph{yīguì}
  \item \emph{Say it:} \emph{yīguì}, once more
  \item \emph{Say it:} \emph{nǐ de dàyī zài yīguì lǐ} — Your coat is in the wardrobe.
  \item \emph{Recall:} say the Mandarin for to hang out to dry, then the Mandarin for to iron, then say \emph{yīguì} again
\end{itemize}

\begin{tcolorbox}[breakable,colback=teal!4,colframe=teal!35!black,title={Before you move on}]
What does yīguì mean? (\textquotedblleft{}Wardrobe\textquotedblright{}.) Say it once more.
\end{tcolorbox}
